\documentclass[11pt,a4paper]{article}

% ── Codificación y fuentes ────────────────────────────────────────────────────
\usepackage[utf8]{inputenc}
\usepackage[T1]{fontenc}
\usepackage{lmodern}
\usepackage[spanish]{babel}

% ── Geometría y espaciado ─────────────────────────────────────────────────────
\usepackage[top=2.5cm, bottom=2.5cm, left=2.8cm, right=2.8cm]{geometry}
\usepackage{setspace}
\setstretch{1.15}
\usepackage{parskip}

% ── Tablas ────────────────────────────────────────────────────────────────────
\usepackage{booktabs}
\usepackage{tabularx}
\usepackage{array}
\usepackage{longtable}
\usepackage{enumitem}

% ── Colores y cajas ───────────────────────────────────────────────────────────
\usepackage[dvipsnames]{xcolor}
\usepackage{tcolorbox}
\tcbuselibrary{skins, breakable}

% ── Cabeceras y pies de página ────────────────────────────────────────────────
\usepackage{fancyhdr}
\usepackage{lastpage}
\pagestyle{fancy}
\fancyhf{}
\fancyhead[L]{\small\textbf{Informe de Evaluación} --- Física para Bioanalistas}
\fancyhead[R]{\small barbara rodriguez 32640760}
\fancyfoot[C]{\small Página \thepage\ de \pageref{LastPage}}
\renewcommand{\headrulewidth}{0.4pt}

% ── Hiperenlaces ──────────────────────────────────────────────────────────────
\usepackage[colorlinks=true, linkcolor=NavyBlue, urlcolor=NavyBlue]{hyperref}

% ── Paleta de colores personalizados ─────────────────────────────────────────
\definecolor{desempeno_excelente}{HTML}{1A6B2F}
\definecolor{desempeno_bueno}{HTML}{2E5A9C}
\definecolor{desempeno_suficiente}{HTML}{8B6914}
\definecolor{desempeno_deficiente}{HTML}{C0392B}
\definecolor{desempeno_insuficiente}{HTML}{7B241C}
\definecolor{grisFondo}{HTML}{F5F5F5}
\definecolor{azulTitulo}{HTML}{1B2A6B}
\definecolor{verdeFortaleza}{HTML}{1A6B2F}
\definecolor{naranjaMejora}{HTML}{A04000}

% ── Estilos de cajas ─────────────────────────────────────────────────────────
\tcbset{
  cajaTitulo/.style={
    enhanced, breakable,
    colback=azulTitulo!8, colframe=azulTitulo,
    fonttitle=\bfseries\large, coltitle=white,
    attach boxed title to top left={yshift=-2mm, xshift=4mm},
    boxed title style={colback=azulTitulo},
    top=6pt, bottom=6pt, left=8pt, right=8pt,
  },
  cajaFortaleza/.style={
    enhanced, breakable,
    colback=verdeFortaleza!6, colframe=verdeFortaleza!60,
    fonttitle=\bfseries, coltitle=verdeFortaleza,
    top=4pt, bottom=4pt, left=8pt, right=8pt,
  },
  cajaMejora/.style={
    enhanced, breakable,
    colback=naranjaMejora!6, colframe=naranjaMejora!60,
    fonttitle=\bfseries, coltitle=naranjaMejora,
    top=4pt, bottom=4pt, left=8pt, right=8pt,
  },
  cajaRecomendacion/.style={
    enhanced, breakable,
    colback=grisFondo, colframe=gray!50,
    fonttitle=\bfseries, coltitle=black,
    top=4pt, bottom=4pt, left=8pt, right=8pt,
  },
}

% ─────────────────────────────────────────────────────────────────────────────
\begin{document}

% ══ PORTADA ══════════════════════════════════════════════════════════════════
\begin{center}
  {\Large\bfseries\color{azulTitulo} Universidad de Oriente/Núcleo Sucre --- Física para Bioanalistas}\\[4pt]
  {\large Informe de Evaluación Preliminar}\\[12pt]
  \rule{\linewidth}{1.2pt}\\[6pt]
  {\LARGE\bfseries barbara rodriguez 32640760}\\[4pt]
  \rule{\linewidth}{0.6pt}
\end{center}

% ══ FICHA TÉCNICA ════════════════════════════════════════════════════════════
\vspace{4pt}
\begin{tcolorbox}[cajaTitulo, title=Datos del informe]
\begin{tabular}{@{}llll@{}}
  \textbf{Estudiante:}  & barbara rodriguez 32640760  &
  \textbf{Fecha:}       & 2026-06-25 \\[4pt]
  \textbf{Archivo PDF:} & \texttt{barbara\_rodriguez\_32640760.pdf}  &
  \textbf{Páginas:}     & 4 \\
\end{tabular}
\end{tcolorbox}

% ══ RESULTADO GLOBAL ═════════════════════════════════════════════════════════
\vspace{6pt}
\begin{center}
  \begin{tcolorbox}[
    enhanced,
    colback=white, colframe=desempeno_excelente,
    width=0.62\linewidth,
    boxrule=2pt,
    halign=center, valign=center,
    top=8pt, bottom=8pt,
  ]
    \centering
    {\large\bfseries Puntaje sugerido}\\[6pt]
    {\Huge\bfseries\color{desempeno_excelente} 9.8/10}\\[6pt]
    {\large Nivel de desempeño: \textbf{\color{desempeno_excelente} Excelente}}
  \end{tcolorbox}
\end{center}

% ══ RESUMEN GENERAL ══════════════════════════════════════════════════════════
\section*{Resumen general}
El estudiante demuestra un excelente dominio de los principios de la física de fluidos aplicados al bioanálisis. Todas las preguntas fueron abordadas con una comprensión conceptual sólida y un desarrollo procedimental mayormente correcto. La organización y claridad de las respuestas son destacables, con un uso apropiado de fórmulas y unidades. Se detectó un error menor de transcripción en un cálculo intermedio, pero el resultado final fue correcto, lo que sugiere una buena comprensión general.

% ══ FORTALEZAS Y ÁREAS DE MEJORA ════════════════════════════════════════════
\vspace{4pt}
\begin{minipage}[t]{0.48\linewidth}
  \begin{tcolorbox}[cajaFortaleza, title={Fortalezas}]
\begin{itemize}[leftmargin=1.5em, itemsep=2pt]
  \item Fuerte comprensión conceptual de los principios de la hidrostática, el principio de Arquímedes, la ecuación de continuidad, la ecuación de Bernoulli y la ley de Poiseuille.
  \item Habilidad para aplicar correctamente las fórmulas y realizar los despejes necesarios.
  \item Consistencia y corrección en el manejo de unidades y conversiones (ej. mm a m).
  \item Presentación clara y organizada de los datos, procedimientos y resultados.
  \item Precisión en la mayoría de los cálculos numéricos.
\end{itemize}
  \end{tcolorbox}
\end{minipage}
\hfill
\begin{minipage}[t]{0.48\linewidth}
  \begin{tcolorbox}[cajaMejora, title={Areas de mejora}]
\begin{itemize}[leftmargin=1.5em, itemsep=2pt]
  \item Revisión cuidadosa de los pasos intermedios de los cálculos para evitar errores de transcripción o presentación, incluso si el resultado final es correcto.
\end{itemize}
  \end{tcolorbox}
\end{minipage}

% ══ OBSERVACIONES POR PREGUNTA ═══════════════════════════════════════════════
\section*{Observaciones por pregunta}

\begin{longtable}{>{\bfseries}p{2.8cm} p{1.6cm} p{7.0cm} p{2.8cm}}
  \toprule
  \textbf{Pregunta} & \textbf{Puntaje} & \textbf{Evaluación} & \textbf{Errores detectados} \\
  \midrule
  \endhead
  \bottomrule
  \endfoot
    \textbf{Pregunta 1: Presión Hidrostática y Venosa} & 2.0/2.0 & La respuesta es completamente correcta. El estudiante aplicó adecuadamente el concepto de presión hidrostática para determinar la altura mínima requerida. Los cálculos son precisos y las unidades son consistentes. La conclusión es clara y pertinente. & \textit{---} \\
    \midrule
    \textbf{Pregunta 2: Principio de Arquímedes} & 2.0/2.0 & La solución es impecable. El estudiante calculó correctamente la fuerza de empuje, el volumen del tejido y su densidad, aplicando de forma precisa el principio de Arquímedes y la definición de densidad. Los cálculos son exactos y la presentación es muy clara. & \textit{---} \\
    \midrule
    \textbf{Pregunta 3: Hemodinámica (Continuidad y Bernoulli)} & 2.0/2.0 & La respuesta es totalmente correcta. El estudiante aplicó de manera excelente la ecuación de continuidad para hallar la velocidad en la zona estrecha y luego la ecuación de Bernoulli para calcular la presión en dicha zona. Todos los pasos son lógicos, los cálculos son correctos y el manejo de unidades es apropiado. & \textit{---} \\
    \midrule
    \textbf{Pregunta 4: Ley de Poiseuille (Análisis Proporcional)} & 2.0/2.0 & La solución es correcta y bien fundamentada. El estudiante identificó correctamente la relación de proporcionalidad de la Ley de Poiseuille con el radio (Q $\alpha$ r$^{4}$) y calculó el porcentaje de caudal mantenido de forma precisa. La interpretación del resultado es adecuada. & \textit{---} \\
    \midrule
    \textbf{Pregunta 5: Flujo Viscoso y Resistencia (Poiseuille)} & 1.8/2.0 & La aplicación de la Ley de Poiseuille para calcular la diferencia de presión es conceptualmente correcta. El despeje de la fórmula y el cálculo de r$^{4}$ son acertados. Sin embargo, se observa un error de transcripción en el cálculo intermedio del numerador (se escribe 7.6 x 10$^{-}$$^{1}$$^{0}$ en lugar de 1.6 x 10$^{-}$$^{1}$$^{0}$). A pesar de este error en la presentación del paso, el resultado final de $\Delta$P (1989 Pa) es correcto, lo que sugiere que el cálculo final se realizó con el valor correcto o que hubo una compensación de errores. Se penaliza ligeramente por la falta de claridad en la presentación del cálculo intermedio. & \textit{Error de transcripción en el cálculo intermedio del numerador (8 * $\eta$ * L * Q)} \\
    \midrule
\end{longtable}

% ══ ERRORES DETECTADOS ═══════════════════════════════════════════════════════
\section*{Análisis de errores}

\subsection*{Errores conceptuales}
\textit{Ninguno identificado.}

\subsection*{Errores procedimentales}
\begin{itemize}[leftmargin=1.5em, itemsep=2pt]
  \item Error de transcripción en un paso intermedio de cálculo en la Pregunta 5.
\end{itemize}

% ══ RECOMENDACIONES AL ESTUDIANTE ════════════════════════════════════════════
\section*{Retroalimentación para el estudiante}
\begin{tcolorbox}[cajaRecomendacion, title={Dirigido a: barbara rodriguez 32640760}]
Bárbara, tu desempeño en este examen es sobresaliente. Demuestras una comprensión profunda de los principios de la física de fluidos y su aplicación en el contexto del bioanálisis. Tus soluciones son claras, organizadas y en su mayoría, muy precisas. La única área de mejora que te sugiero es ser aún más meticulosa al escribir los pasos intermedios de tus cálculos. En la Pregunta 5, aunque tu resultado final fue correcto, hubo un error de transcripción en un valor intermedio. Prestar atención a estos detalles asegura que tu trabajo sea impecable y fácil de seguir. ¡Sigue así, tu nivel es excelente!
\end{tcolorbox}

% ══ NOTA PARA EL PROFESOR ════════════════════════════════════════════════════
\section*{Nota para el profesor}
\begin{tcolorbox}[
  enhanced, breakable,
  colback=yellow!8, colframe=yellow!60!black,
  fonttitle=\bfseries, coltitle=black,
  title={Observaciones para revisión manual},
  top=4pt, bottom=4pt, left=8pt, right=8pt,
]
El estudiante Bárbara Rodríguez ha presentado un examen de muy alto nivel. La comprensión conceptual es sólida en todas las preguntas y la aplicación de los principios es correcta. El único punto a destacar es un error de transcripción en un cálculo intermedio de la Pregunta 5; sin embargo, el resultado final de esa pregunta es correcto, lo que sugiere que el estudiante realizó el cálculo correctamente en su calculadora pero lo transcribió mal. Considero que la penalización aplicada es justa dada la calidad general del trabajo.
\end{tcolorbox}

% ══ PIE DE INFORME ════════════════════════════════════════════════════════════
\vspace{12pt}
\begin{center}
  \footnotesize\color{gray}
  Informe generado automáticamente por el sistema de evaluación con IA (gemini-2.5-flash) y soporte LaTex. \\
  Este documento es una evaluación \textbf{preliminar} para ser revisado por el profesor antes de ser entregado al 
  estudiante. \\
  Generado el 2026-06-25 \textbullet\ BIOANALISIS Sistema Académico de Evaluación.
  \end{center}

\end{document}
